\chapter{Technical interfaces}
\label{mk-app-technical}

The main chapters state the mathematical route.  These interface nodes record
where the detailed formalization companion enters without exposing every
transport lemma in the book's dependency graph.

\begin{proposition}[Cech computation interface]
  \label{mk-app-cech-interface}
  \dcref{III.1,III.12,custom}
  \uses{mk-prop-cech-computes,mk-lem-cech-flat-cohomology,mk-thm-cohomology-base-change}
  \cite{stacks-cohomology-schemes,stacks-project}
  On a finite affine cover of a relative curve, the Cech complex computes
  coherent cohomology and is preserved by field extension.  Consequently its
  kernels and cokernels give the base-change maps used in the Picard tangent
  calculation.
\end{proposition}

\begin{proof}
  Affine acyclicity identifies derived cohomology with the Cech complex.
  Restriction maps and localization commute with tensoring by a field, and
  flatness carries this identification through kernels and cokernels.
\end{proof}

\begin{proposition}[Picard parameter-space interface]
  \label{mk-app-picard-parameter-interface}
  \dcref{4.8,4.9,custom}
  \uses{mk-thm-hilbert-quot,mk-lem-flattening-stratification,mk-thm-picard-representability}
  \cite{kleiman-picard,nitsure-hilbert-quot}
  The bounded Quot and flattening constructions produce a scheme representing
  the relative Picard functor after imposing the line-bundle and descent
  conditions.
\end{proposition}

\begin{proof}
  A sufficiently positive twist embeds the relevant quotients in a bounded
  parameter space.  Flattening isolates the constant Hilbert polynomial locus;
  the invertible locus is open, and the residual base-line-bundle ambiguity is
  removed by the descent quotient.
\end{proof}

\begin{proposition}[Identity-component interface]
  \label{mk-app-identity-interface}
  \dcref{5.1,5.4,5.11,custom}
  \uses{mk-thm-pic0-abelian,mk-prop-genus-tangent,mk-thm-cohomology-base-change}
  \cite{kleiman-picard}
  The identity component of the represented Picard scheme is smooth and
  proper, and its dimension is the cohomological genus.
\end{proposition}

\begin{proof}
  Large-degree Abel maps prove properness of a degree component.  Translation
  identifies it with the identity component, while the tangent-space
  calculation gives dimension \(\dim_k H^1(C,\mathcal O_C)=g(C)\); vanishing in
  degree two removes the infinitesimal obstruction.
\end{proof}

\begin{theorem}[Rigidity and extension interface]
  \label{mk-app-rigidity-interface}
  \dcref{I.2,3.1,3.2,custom}
  \uses{mk-thm-rigidity,mk-thm-rational-extension}
  \cite{milne-av,mumford-av}
  A pointed morphism from a product with a proper connected factor is rigid,
  and a rational map from a nonsingular variety to an abelian variety extends
  uniquely.  These are the two geometric inputs for the Albanese property.
\end{theorem}

\begin{proof}
  Rigidity makes the difference of two candidate maps constant on the proper
  factor.  For a rational map, properness gives a graph closure and the
  valuative criterion gives uniqueness at codimension one; rigidity removes
  the remaining indeterminacy.
\end{proof}

\begin{proposition}[Symmetric-power interface]
  \label{mk-app-symmetric-power-interface}
  \dcref{III.6,custom}
  \uses{mk-thm-symmetric-powers,mk-thm-abel-jacobi-universal}
  \cite{milne-av}
  The symmetric powers of the curve carry the descended Abel map, and their
  complete-linear-system fibres are exactly what makes maps to abelian
  varieties factor through the Jacobian.
\end{proposition}

\begin{proof}
  The ordered sum is invariant under permutation and therefore descends.
  Equal line-bundle classes differ by a principal divisor, so each fibre is a
  complete linear system.  Rigidity makes a map constant on that fibre, hence
  it factors through the Abel map.
\end{proof}

\begin{theorem}[Jacobian interface]
  \label{mk-app-jacobian-interface}
  \dcref{III.6,5.13,5.14,custom}
  \uses{mk-thm-jacobian-package}
  \cite{milne-av,kleiman-picard}
  The final Jacobian package is obtained by composing the Picard
  representability, identity-component, rigidity, and Abel--Jacobi interfaces.
\end{theorem}

\begin{proof}
  The preceding interfaces identify \(\operatorname{Jac}(C)\) with the
  abelian variety \(\operatorname{Pic}^0_{C/k}\), compute its dimension, and
  supply the universal factorisation and base-change compatibilities.
\end{proof}
